\documentclass[11pt]{article}
\usepackage[margin=1in]{geometry}
\usepackage{amsmath,amssymb,amsthm,mathtools}
\usepackage{array,booktabs,longtable}
\usepackage[hidelinks]{hyperref}

\newtheorem{theorem}{Theorem}[section]
\newtheorem{lemma}[theorem]{Lemma}
\newtheorem{proposition}[theorem]{Proposition}
\newtheorem{corollary}[theorem]{Corollary}
\theoremstyle{definition}
\newtheorem{definition}[theorem]{Definition}
\theoremstyle{remark}
\newtheorem{remark}[theorem]{Remark}

\newcommand{\spn}{s^+}
\newcommand{\smn}{s^-}
\newcommand{\betarank}{\beta}
\newcommand{\one}{\mathbf 1}

\title{Positive Square Energy of Every Hexacyclic Graph}
\author{AI-generated manuscript; no human authorship is claimed}
\date{12 August 2026}

\begin{document}
\maketitle

\begin{abstract}
Let $G$ be a finite simple connected graph with $n$ vertices and $n+5$ edges.
We prove that the sum $\spn(G)$ of the squares of its positive adjacency
eigenvalues is at least $n$.  Cyclomatic-rank additivity splits the proof into
graphs with at least two positive-rank cyclic blocks and graphs with exactly
one.  The first branch is covered by an owner-exact multiblock ledger.  In the
second branch, suppression gives one of $1198$ rank-six multigraph kernels,
with order distribution $1,4,26,84,216,314,325,162,66$ on two through ten
branch vertices.  Exact finite owners cover the complete kernel partition;
exact path elimination promotes their physical parity frontiers to arbitrary
simple subdivisions, and one-vertex additivity includes arbitrary rooted
trees.  A fail-closed all-connected master authenticates both branches, checks
that their scopes are disjoint and exhaustive, and preserves the nonstrict
conclusion.  The computational layer uses committed exact-execution receipts;
their byte authentication is distinguished explicitly from a fresh independent
replay.
\end{abstract}

\section{Statement and scope}

All graphs in the spectral statements are finite, simple, and undirected.  For
the adjacency eigenvalues $\lambda_1,\ldots,\lambda_n$ of $G$, put
\[
 \spn(G)=\sum_{\lambda_i>0}\lambda_i^2,
 \qquad
 \smn(G)=\sum_{\lambda_i<0}\lambda_i^2,
 \qquad
 \betarank(G)=|E(G)|-|V(G)|+1.
\]
A connected graph with $|E(G)|=|V(G)|+5$ has cyclomatic rank six.

Positive and negative square energies were introduced in connection with
spectral coloring bounds by Elphick, Farber, Goldberg, and Wocjan
\cite{EFGW}.  Liu, Tang, and Zhang proved the conjectured general bound
$\min\{\spn(G),\smn(G)\}\ge n-1$ for connected $n$-vertex graphs
\cite{LTZ}.  The conclusion considered here is instead the stronger
edge-surplus conjecture of Akbari, Kumar, Mohar, Pragada, and Zhang: a connected
$n$-vertex graph with at least $n+1$ edges should satisfy $\spn(G)\ge n$
\cite[Conjecture~1.2]{AKMPZ}.  The present result proves its exact hexacyclic
frontier; it does not infer cases with more edges by edge-addition monotonicity.

\begin{theorem}\label{thm:main}
Every finite simple connected graph $G$ with $|V(G)|=n$ and
$|E(G)|=n+5$ satisfies
\[
 \boxed{\spn(G)\ge n}.
\]
\end{theorem}

No strict inequality, equality classification, edge-addition monotonicity,
one-edge-subdivision monotonicity, or result for nonsimple realizations is
claimed.  Human review, independent verification, peer review, and publication
are likewise not claimed.

\section{Analytic foundations}

For a finite graph $H=(V,E)$ define
\begin{equation}\label{eq:kappa}
 \kappa(H)=\min_{\substack{R\succeq0\\R_{vv}=1\ (v\in V)}}
 \sum_{uv\in E}\frac{2}{1-R_{uv}},
\end{equation}
where a summand is $+\infty$ if $R_{uv}=1$.

\begin{lemma}[DNN witness implication]\label{lem:dnn}
For every finite simple graph $H$,
\[
 \smn(H)\le\kappa(H).
\]
Consequently, any feasible correlation matrix with objective at most $b$
proves $\spn(H)\ge2|E(H)|-b$.
\end{lemma}
\begin{proof}
Write the adjacency matrix as $A=P-N$, where $P,N\succeq0$ are its positive
and negative spectral parts, and put $M=N\circ N$.  The Schur product theorem
gives $M\succeq0$ and $M\ge0$ entrywise.  Moreover,
\[
 \langle J,M\rangle=\operatorname{tr}N^2=\smn(H),
 \qquad
 \frac{\smn(H)}2=-\sum_{uv\in E(H)}N_{uv}
 \le\sum_{uv\in E(H)}\sqrt{M_{uv}}.
\]
For every feasible $R$,
\[
 \langle J,M\rangle-2\sum_{uv\in E(H)}(1-R_{uv})M_{uv}
 =\langle R,M\rangle+
 2\sum_{\substack{u<v\\uv\notin E(H)}}(1-R_{uv})M_{uv}\ge0.
\]
Weighted Cauchy--Schwarz and cancellation when $\smn(H)>0$ give
\[
 \smn(H)\le\sum_{uv\in E(H)}\frac{2}{1-R_{uv}}.
\]
The zero and infinite-objective cases are immediate.  Minimize over $R$ and
use $\spn(H)+\smn(H)=\operatorname{tr}A^2=2|E(H)|$.
\end{proof}

\begin{lemma}[Exact path cost]\label{lem:path}
If a kernel edge with endpoint correlation $r$ is replaced by a path of length
$q\ge1$, its least DNN contribution in excess of its $q$ edges is
\[
 f_q(r)=q\tan^2\!\left(
 \frac{\arccos((-1)^q r)}{2q}\right).
\]
For fixed $r$ and parity, $f_{q+2}(r)\le f_q(r)$, with endpoint values
understood in the extended sense.
\end{lemma}
\begin{proof}
Represent the endpoints by unit vectors and alternately negate vectors along
the path.  The transformed endpoint angle is
$\gamma=\arccos((-1)^q r)$.  The spherical triangle inequality and convexity
of $\sec^2(x/2)$ make $q$ equal angles optimal, giving
$q\sec^2(\gamma/(2q))=q+f_q(r)$.  With $z=\gamma/(2q)$,
\[
 \frac{d}{dq}\left[q\tan^2\!\left(\frac\gamma{2q}\right)\right]
 =\tan z\sec^2z\,(\sin z\cos z-2z)\le0.
\]
One-sided limits handle endpoint cases.
\end{proof}

\begin{lemma}[One-vertex additivity]\label{lem:additivity}
If $H=H_1\vee H_2$ is a one-vertex sum, then
$\kappa(H)=\kappa(H_1)+\kappa(H_2)$.  In particular,
$\kappa(T)=|E(T)|$ for every finite tree $T$.
\end{lemma}
\begin{proof}
Restriction gives one inequality.  For the other, identify the common unit
vector in minimizing Gram representations and place the two perpendicular
components in mutually orthogonal subspaces.  The objectives add.  Every tree
edge costs at least one, and opposite unit vectors on the two bipartition
classes attain one on each edge.
\end{proof}

\section{Block decomposition}

We use the standard block convention: the blocks of a connected graph are its
maximal $2$-connected subgraphs together with its bridge copies of $K_2$.  A
\emph{positive-rank cyclic block} is a block $B$ with $\betarank(B)>0$.
Thus ``a unique positive-rank cyclic block'' does not mean ``one nontrivial
block'': attached trees contribute nontrivial $K_2$ blocks.

\begin{lemma}[Cyclic-block rank additivity]\label{lem:block-rank}
The cyclomatic rank of a connected graph is the sum of the cyclomatic ranks of
its positive-rank cyclic blocks.
\end{lemma}
\begin{proof}
Following the block--cut tree, adjoining a block $B$ at one cut vertex adds
$|V(B)|-1$ vertices and $|E(B)|$ edges, and hence adds
$|E(B)|-|V(B)|+1=\betarank(B)$ to the rank.  Bridge blocks contribute zero.
\end{proof}

Thus the positive cyclic-block ranks of a rank-six graph form exactly one of
\begin{gather}\label{eq:partitions}
 1^6,\quad2+1^4,\quad2+2+1+1,\quad2+2+2,\quad3+1+1+1,\quad
 3+2+1,\quad3+3,\nonumber\\
 4+1+1,\quad4+2,\quad5+1,\quad6.
\end{gather}
The first ten alternatives are the multiblock branch; the final alternative has one
positive-rank cyclic block.  This is a structural dichotomy only and does not
by itself prove either branch.

\begin{lemma}[Form of the unique-block branch]\label{lem:unique-block-form}
If a connected graph has a unique positive-rank cyclic block $B$, then every
other block is a bridge.  The graph is obtained from $B$ by one-vertex-summing
finite rooted trees at vertices of $B$.
\end{lemma}
\begin{proof}
Every other block has rank zero.  Under the stated block convention, a
rank-zero block of a simple graph is a bridge $K_2$.  Delete the node for $B$
from the block--cut tree.  Each remaining component is a tree of bridge blocks
attached through exactly one cut vertex of $B$; two attachments to $B$ would
create a cycle in the block--cut tree.  Expanding its bridge blocks gives the
asserted rooted tree.  After suppression is reversed, an attachment vertex may
be an internal vertex of a replacement path.
\end{proof}

\section{Multiblock branch}

\begin{proposition}[Multiblock branch]\label{prop:multiblock}
Every finite simple connected rank-six graph with at least two positive-rank
cyclic blocks satisfies $\spn(G)\ge|V(G)|$.
\end{proposition}

The cited owner-exact multiblock audit \cite{MultiblockAudit} proves the
proposition and supplies its complete repository ledger.  Its scope consists of
the first ten rows of \eqref{eq:partitions}.  It separates direct DNN rows,
structural pre-sieve rows, and the packet residuals A--I.  Shared cuts,
repeated and nested owners, positive bridge connectors, and arbitrary rooted
branches are assigned before credit is charged.  The direct
rank-five-equality-plus-triangle boundary is nonstrict and is not promoted to a
strict claim.  This skeleton cites that internal exposition rather than
duplicating its long owner-by-owner argument.  A successful ledger replay
checks the encoded predicates; it does not by itself establish that every
encoded structural lemma is mathematically sound.

\section{Unique positive-rank block}

\begin{lemma}[Rank-six suppression]\label{lem:suppression}
Let $B$ be a finite simple $2$-connected graph with $\betarank(B)=6$.
Suppressing every degree-two vertex produces a loopless $2$-connected
multigraph $K$ of minimum degree at least three such that
\[
 |E(K)|=|V(K)|+5,
 \qquad 2\le|V(K)|\le10.
\]
Conversely, $B$ is a simple realization obtained by replacing kernel edges by
positive-length, internally disjoint paths meeting only at prescribed branch
endpoints.
\end{lemma}
\begin{proof}
Suppression preserves $|E|-|V|$, connectedness, and rank.  A loop would come
from a cycle meeting the remainder only at one branch vertex, contradicting
$2$-connectivity.  If $k=|V(K)|$, minimum degree three and
$|E(K)|=k+5$ give $3k\le2k+10$, so $k\le10$; loopless $2$-connectivity gives
$k\ge2$.  Reversing suppression gives the stated path model.  Simplicity is a
separate required condition, especially in parallel classes.
\end{proof}

\begin{proposition}[Exact kernel census]\label{prop:census}
The rank-six suppressed kernels have the following repository-checked
distribution.
\[
\begin{array}{c|rrrrrrrrr}
 |V(K)|&2&3&4&5&6&7&8&9&10\\ \hline
 \#K&1&4&26&84&216&314&325&162&66
\end{array}
\]
There are $1198$ isomorphism classes in total.
\end{proposition}

The exact integer census regenerates multiplicity vectors, rejects supports
with a cut vertex, canonicalizes under isomorphism, and compares the complete
sorted output with \path{research/fixtures/rank-six-kernels.json}.  No
floating-point or spectral calculation enters this census.  Its reproducible
command is listed in Appendix~\ref{app:repro}.

\section{Finite-owner interface}

Here an \emph{exact DNN certificate of excess at most five} means an explicit
finite record which reconstructs a positive-semidefinite Gram matrix of unit
vectors on the realized core, verifies every shared-endpoint identity and every
edge denominator $1-R_{uv}>0$, and proves by exact integer, rational, or
symbolically justified arithmetic that
\[
 \sum_{uv\in E}\frac{2}{1-R_{uv}}\le |E|+5.
\]
A structural owner is admissible only if it has a separately stated proof of
the same final spectral conclusion for its exact scope; it is not a DNN
certificate merely because a verifier assigns it a target key.  Numerical
optimization output, a symbolic tag without checked identities, or a stored
decimal approximation is not an exact certificate.

For each parallel class of multiplicity $m$ with $o$ odd replacement paths,
the canonical simple lengths are
\[
 c(m,o)=
 \begin{cases}
 (2,\ldots,2),&o=0,\\
 (1,\underbrace{3,\ldots,3}_{o-1},
       \underbrace{2,\ldots,2}_{m-o}),&o>0.
 \end{cases}
\]
Simplicity permits at most one unit path in a parallel class.  After permuting
parallel physical edges, every realization vector $\ell$ in the parity row
satisfies $c_i\le\ell_i$ and $\ell_i-c_i\in2\mathbb Z_{\ge0}$.

\begin{proposition}[Exact-owner lift]
\label{prop:owner-interface}
Suppose every physical parity orbit of every kernel in
Proposition~\ref{prop:census} is owned either by an exact coarse DNN certificate
at its canonical vector whose fixed branch Gram remains within excess five
after exact path elimination under every same-parity coordinate lengthening,
or by exact DNN certificates (or explicitly typed structural owners proving
the same conclusion) for every target
\[
 \mathcal F(c)=\{c\}\cup\{c+2\mathbf e_i:1\le i\le|E(K)|\}.
\]
For a structural owner, assume additionally an exact all-length and rooted-tree
lift covering the descendants for which it claims ownership;
Lemma~\ref{lem:path} may be used only for a DNN owner with a retained branch
Gram.  Assume that regenerated target keys equal the disjoint union of checked
final-owner keys, with no missing, duplicate, or extra keys.  Then every simple
subdivision of every rank-six kernel, with arbitrary
finite rooted trees attached at branch or internal subdivision vertices,
satisfies $\spn(G)\ge|V(G)|$.
\end{proposition}
\begin{proof}
If $\ell=c$, use the canonical target.  Otherwise choose $i$ with
$\ell_i\ge c_i+2$ and start from the owner of $c+2\mathbf e_i$.  For a DNN
owner, retain its branch vectors, replace the internal vectors independently
on every lengthened path, and apply Lemma~\ref{lem:path}; the excess cannot
increase.  A typed structural owner instead uses the all-length and rooted-tree
implication required in the premise.  Thus in the DNN case the subdivided core
$B$, with $L$ edges, has $\kappa(B)\le L+5$.

In that DNN case, if attached rooted trees contain $t$ edges,
Lemma~\ref{lem:additivity} gives
$\kappa(G)\le L+5+t$.  Rank six gives $|V(G)|=L-5+t$, and
Lemma~\ref{lem:dnn} yields
\[
 \spn(G)\ge2(L+t)-(L+5+t)=L-5+t=|V(G)|.
\]
In the structural case, the all-length and rooted-tree conclusion stipulated
in the premise gives the same inequality directly.
Different frontier targets may use different Gram matrices.  No parity-changing
subdivision or spectral monotonicity is used.
\end{proof}

The orders-$2$--$10$ theorem master verifies this premise over the exact
partition $K1$--$K645$, $K646$--$K970$, $K971$--$K1132$, and
$K1133$--$K1198$.  Its four owners cover orders $2$--$7$, $8$, $9$, and $10$,
respectively, and its regenerated counts are
$1,4,26,84,216,314,325,162,66$.  The master authenticates committed canonical
exact-execution records for every child, checks the analytic lift separately,
and rejects missing, overlapping, widened, or status-only owners.

\section{Assembly}

\begin{proof}[Proof of Theorem~\ref{thm:main}]
By Lemma~\ref{lem:block-rank}, the positive block ranks have exactly the eleven
forms in \eqref{eq:partitions}.  If there are at least two positive-rank cyclic
blocks, apply Proposition~\ref{prop:multiblock}.  Otherwise there is one
rank-six cyclic block.  Lemma~\ref{lem:unique-block-form},
Lemma~\ref{lem:suppression},
Proposition~\ref{prop:census}, the verified orders-$2$--$10$ owner partition,
and Proposition~\ref{prop:owner-interface} prove the second branch.  The
all-connected master checks that the two branch predicates are disjoint and
exhaustive and that both owners state the same nonstrict conclusion.  This
proves the theorem.
\end{proof}

\appendix
\section{Reproducibility ledger}\label{app:repro}

\subsection{Commands for present repository artifacts}

From the repository root, the kernel census and multiblock ledger can be
replayed in normal and optimized modes:
\begin{verbatim}
python3 research/rank-six-kernel-census-verifier.py
python3 -O research/rank-six-kernel-census-verifier.py
python3 research/hexacyclic-multiblock-ledger-verifier.py
python3 -O research/hexacyclic-multiblock-ledger-verifier.py
\end{verbatim}
The complete single-block and all-connected theorem owners are audited with
\begin{verbatim}
python3 research/rank-six-order2-10-master-verifier.py
python3 -O research/rank-six-order2-10-master-verifier.py
python3 research/hexacyclic-all-connected-master-verifier.py
python3 -O research/hexacyclic-all-connected-master-verifier.py
\end{verbatim}

\subsection{Frozen theorem-owner identities}

The theorem conversion is based on committed revision
\texttt{3ac4e68d2ef30b28925d5c30563b5e66ebfae4c0}.  The final source and
canonical-output SHA-256 values are:
\begin{center}
\small
\begin{tabular}{p{0.37\linewidth}p{0.26\linewidth}p{0.26\linewidth}}
\toprule
owner&source SHA-256&output SHA-256\\\midrule
order nine&\texttt{d11b6da30ccff5905}&\texttt{dd6a5a1e9b905955}\\
order ten&\texttt{4cb22611448755d1}&\texttt{767bbf4dd142afd7}\\
orders $2$--$10$ master&\texttt{21aabed4eee24f36}&\texttt{e32149c665638e1d}\\
all-connected master&\texttt{cd22798e32c1f51b}&\texttt{de4f7ff814530bd0}\\
\bottomrule
\end{tabular}
\end{center}
The displayed values are abbreviated for typesetting.  Their complete values,
which are the values checked by the verifiers and the reproduction record, are
respectively
\begin{verbatim}
d11b6da30ccff5905cf851d33fc6a5d0cfc6650f3e8317224ddc0d8a2e4cd3d9
dd6a5a1e9b905955affde5fa64e0b8187dc39ad35c02ab007172da3358dff4f4
4cb22611448755d1eea984271251b542ab96f519cc052f3aeda8d56c04c61819
767bbf4dd142afd7eb5a45dffeb738734948412ca53358c40d3655b28403725e
21aabed4eee24f36dc9ddf0460423d0fdd6d20a56e80a60ffed69c796c163bf1
e32149c665638e1d9cd08f8050cd4a7c7cae9f76c51b49f9cda3e0cfb2645834
cd22798e32c1f51bf5a1646eb063215731c30ee14b4570c60b9fb1faf8584379
de4f7ff814530bd0a69c4d6a1407cf5583bff8dee3aeab41600b3e971afd127e
\end{verbatim}
The committed child-execution evidence has SHA-256
\begin{verbatim}
bbbb8eff45f8009dfb49e926c4952ae0b231ed6e02e708bddfb13b4aac17bb5d
\end{verbatim}
Normal and optimized execution on Python 3.12.3 produced byte-identical
orders-$2$--$10$ and all-connected outputs with the complete output digests
listed above.

\subsection{Build}

This manuscript is compiled from the repository root with
\begin{verbatim}
bash scripts/build-paper.sh all-hexacyclic-graphs
\end{verbatim}
Build success checks only the LaTeX artifact; the theorem evidence is supplied
by the owner chain above, not by typesetting.

\section{Computational certificate specification}
\label{app:certificate}

This appendix specifies the formats and acceptance algorithm used by the
order-nine and order-ten finite owners.  The final theorem owners require full
manifest coverage and exact final ownership; an \texttt{unresolved} record is
a valid encoding but is never an accepting proof record.

\subsection{Common conventions and target keys}

All integers below are mathematical integers, not machine words.  A
\emph{uvarint} writes a nonnegative integer in base $128$, least significant
group first, with bit $7$ set exactly when another group follows.  A signed
integer $z$ is first mapped to $2z$ for $z\ge0$ and to $-2z-1$ for $z<0$, then
written as a uvarint.  Multi-byte integers are therefore independent of host
endianness.  The final decoders reject a uvarint with a redundant final zero
group by requiring decode--re-encode byte equality, so two encodings of one
certificate cannot be accepted.

The frozen kernel fixture, in its listed order, regenerates the residual source
stream.  A source index $j$ denotes row $j$ of that stream (zero based), not a
kernel identifier stored in a pack.  The verifier reconstructs from $j$ the
kernel, its ordered physical paths, the parity row, canonical lengths, and the
$p+1$ frontier labels
\[
  C, X_0,\ldots,X_{p-1},
  \qquad p=14\text{ for order nine},\quad p=15\text{ for order ten}.
\]
Thus an exact target key is the pair $(j,C)$ or $(j,X_i)$.  A pack owns a
half-open source interval $[a,a+r)$, encoded by uvarints $a,r$.  Keys are never
accepted from asserted counts: the master regenerates them and compares sets.

Every stored stereographic parameter vector has dimension $d-1$, where
$d\in\{9,10\}$.  One positive common denominator $D$ precedes a witness and
each coordinate is a signed-varint numerator $n$, denoting $n/D$.  Vector and
path counts are implicit in the regenerated source row; no untrusted dimension
field controls allocation or verification.

\subsection{The \texttt{R9G1} wire format}

After XZ decompression, an order-nine pack is the following byte concatenation:
\begin{verbatim}
"R9G1" | fixture_sha256[32] | uvarint(start) | uvarint(count)
       | record[0] | ... | record[count-1].
\end{verbatim}
The 32-byte digest is binary, not 64 hexadecimal characters.  Here $d=9$,
$p=14$, and a target bitmap has 15 significant bits.  Record tags are
\begin{center}
\begin{tabular}{ccl}
\toprule
byte&name&body and meaning\\
\midrule
0&\texttt{unresolved}&empty; owns no target\\
1&\texttt{shared}&one shared rational witness, described below\\
2&\texttt{template}&empty; invokes the exact K971 signed-five-cycle audit\\
3&\texttt{individual}&uvarint bitmap, then one witness per set bit\\
\bottomrule
\end{tabular}
\end{center}
A shared body is $D$, nine branch parameter vectors, then, in physical-path
order, all canonical internal vectors and all length-plus-two internal vectors.
For a canonical path of length $q_i$, these two lists have respectively
$q_i-1$ and $q_i+1$ vectors.  The same branch vectors are used for all 15
frontiers.  An individual body begins with a nonzero bitmap smaller than
$2^{15}$; for each set bit in increasing target order it stores a fresh $D$,
nine branch vectors, and $q_i-1$ internal vectors on each path, except that the
selected $X_i$ path has $q_i+1$.  The template tag is legal only on the
regenerated K971 template row.  The decoder rejects bad magic or fixture hash,
unknown tags, zero denominators, out-of-range intervals or bitmaps, truncation,
template/tag mismatch, and trailing raw bytes.

\subsection{The \texttt{R10G1} wire format}

The order-ten header is identical with magic \texttt{R10G1}; here $d=10$,
$p=15$, and the bitmap has 16 significant bits.  Tags 0--3 have the same
layout, with tag 3 named \texttt{fallback}.  The current version also reserves
payload-free, source-row-dependent exact recognizers:
\begin{center}
\begin{tabular}{ccl}
\toprule
byte&name&required independent audit\\
\midrule
4&\texttt{structural}&the typed structural Gram sieve\\
5&\texttt{atom}&the frozen atom-ledger classification and exact atom audit\\
6&\texttt{balanced}&the balanced rank-one recognizer\\
\bottomrule
\end{tabular}
\end{center}
Tag 2 invokes only the regenerated K1133 signed-five-cycle row.  Empty tags do
not mean ``trust the classification'': their recognizers rebuild the required
combinatorics and exact Gram or structural predicate.  The \texttt{R10G1}
decoder rejects the malformed cases listed for \texttt{R9G1} and requires the
raw input to equal its unique re-encoding.  XZ is a transport envelope only;
acceptance is defined on the decompressed bytes and the exact records they
denote.

\subsection{Hardware-independent exact arithmetic}

For a parameter row $t=(t_1,\ldots,t_{d-1})\in\mathbb Q^{d-1}$ the verifier
constructs, without floating point,
\[
 u(t)=\frac{(1-\langle t,t\rangle,2t_1,\ldots,2t_{d-1})}
             {1+\langle t,t\rangle}.
\]
Python arbitrary-precision integers and normalized \texttt{Fraction} values
implement the field operations.  For consecutive unit vectors $x,y$ in the
alternating-sign path chain it requires $1+\langle x,y\rangle>0$ and adds
\[
             \frac{1-\langle x,y\rangle}{1+\langle x,y\rangle}.
\]
The final comparison is the exact cross-multiplied rational inequality
$\sum\mathrm{cost}\le5$.  Endpoint vectors are reconstructed from the common
branch list, with the terminal branch negated exactly when prescribed by path
parity.  Consequently unit norm, shared endpoints, path widths, denominator
signs, and the budget are rational identities or integer comparisons.  Search
seeds, floating-point optimizer scores, wall-clock times, CPU vector
instructions, byte order, and native integer overflow cannot affect
acceptance.  A conforming alternative verifier may use any big-integer
rational library, but must produce the same accepted key set from the same raw
bytes.  XZ decompression, SHA-256, and Python itself remain implementation
dependencies; this arithmetic claim is not a claim of independent software
verification.

\subsection{Fail-closed verifier algorithm}

For each order, the theorem verifier executes the following algorithm.
Every failure aborts with nonzero status; no \texttt{assert}, status string, or
stored total is an acceptance condition.
\begin{enumerate}
\item Hash and parse the pinned kernel fixture, verifier sources, recognizer
sources, census output, and canonical manifest; reject any unpinned byte
sequence or noncanonical manifest JSON.
\item Regenerate the complete ordered residual source stream and, from it, the
complete expected target-key set.  Compare its count and ordered SHA-256 with
the pinned values.
\item Read manifest chunks in increasing half-open ranges.  Require paths to
remain below the manifest directory, and ranges to be nonempty, contiguous,
nonoverlapping, and exactly $[0,N)$.  Reject repeated paths.
\item For each chunk, check compressed byte count and SHA-256, decompress one XZ
stream, check raw byte count and SHA-256, parse the matching embedded interval,
and reject malformed or noncanonical records and trailing bytes.
\item For every rational record, reconstruct every relevant frontier chain and
perform the exact arithmetic audit above.  For every payload-free record, run
its named exact recognizer.  Map only successfully checked targets to their
typed owner; map no target from an unresolved bit or tag.
\item Regenerate the separately pinned symbolic-owner keys.  Check that each is
proved by its typed recognizer, that rational and final symbolic ownership are
disjoint after the prescribed precedence rule, and that their union equals the
expected target set exactly.  Explicitly reject missing, duplicate, and extra
keys.
\item Emit canonical ASCII JSON containing scopes, dependency digests, mode and
owner counts, ordered key digests, exact partition equality, and the finite
conclusion.  Replay under \texttt{python3} and \texttt{python3 -O};
require zero exits and byte-identical reports.  The orders-$2$--$10$ master must
then validate, rather than copy, these two order conclusions.
\end{enumerate}
Steps 1--4 alone are a digest audit, not certificate verification.  Likewise,
successful verification of one chunk proves only the records in that chunk;
aggregating receipts does not replace a fresh full-manifest exact replay.

\subsection{Segmented reproducibility}

Large packs may be replayed independently by manifest chunk.  From the
repository root, first perform the inexpensive range and byte audit:
\begin{verbatim}
python3 positive-square-energy/experiments/rank6_order9_pack_auditor.py \
  --manifest positive-square-energy/experiments/rank6_order9_search_manifest.json \
  --digest-only
python3 positive-square-energy/experiments/rank6_order10_pack_auditor.py \
  --manifest positive-square-energy/experiments/rank6_order10_search_manifest.json \
  --digest-only
\end{verbatim}
For each manifest chunk index $i=0,1,\ldots,m-1$, use a new output path:
\begin{verbatim}
python3 positive-square-energy/experiments/rank6_order9_pack_auditor.py \
  --manifest positive-square-energy/experiments/rank6_order9_search_manifest.json \
  --chunk-index i --write-chunk-receipt /tmp/r9-replay/chunk-i.json
python3 positive-square-energy/experiments/rank6_order10_pack_auditor.py \
  --manifest positive-square-energy/experiments/rank6_order10_search_manifest.json \
  --chunk-index i --write-chunk-transcript /tmp/r10-replay/chunk-i.json
\end{verbatim}
The literal \texttt{i} is replaced by the decimal index.  The order-ten command
also launches and compares optimized replay; order nine must be run explicitly
under both \texttt{python3} and \texttt{python3 -O} until its auditor provides
the same built-in comparison.  Each segment authenticates the full manifest
and dependency map before auditing its own raw records, so segment boundaries
cannot alter certificate semantics.  A clean full replay requires all chunks,
then the full auditor with no \texttt{--chunk-index}; if coverage is incomplete
or any target remains unowned, the expected result is nonzero status.  Receipts
and transcripts are resumability aids only and are not theorem owners.

\subsection{Final manifest identities}

The frozen kernel-fixture SHA-256 embedded in both streams is
\begin{verbatim}
5a862a0e9ed5dfe91ff6f8491936c8e775eb39b71619df6b8c2a9be2c4643476
\end{verbatim}
The canonical order-nine and order-ten manifest SHA-256 values are
\begin{verbatim}
8aa9d797d9ed786ad438d6fd685e0ec576247b45c17a14749b76c45eebbe9168
5162243e0535ca41f72fa0c7bd27e6ee240d485567e56b8c8dfbcf3a4ecbf3b6
\end{verbatim}
respectively.  Complete leaf source maps, compressed/raw chunk hashes, ranges,
receipt identities, and owner totals remain in those canonical manifests and
their transitive theorem-owner dependencies rather than being duplicated as a
mutable typeset table.

\section{Trusted base and limitations}

The analytic trust base consists of the displayed matrix arguments, exact path
elimination, standard block theory, the suppression reduction, and every
structural lemma imported through the multiblock ledger or a typed finite
owner.  Those imported repository notes are mathematical dependencies, not
external validation.  The computer-assisted trust base additionally includes the
committed verifier and generator sources, canonical kernel fixtures, every
certificate and chunk consumed transitively, decompression and parsing code,
exact integer/rational and validated symbolic semantics, canonical labelling,
the Python runtime, and executing hardware.  Hashes establish byte identity,
not truth; normal/optimized agreement detects only some implementation errors;
hostile mutations are tests, not a completeness proof; and no independent
implementation is claimed.  Numerical optimization may discover witnesses but
is not a proof object.  Rooted-tree attachment means a genuine one-vertex sum;
a connector meeting a core twice is outside that lemma.  Kernel parallel edges
are permitted only before realization; the graph to which the trace identity
is applied remains simple.

\section*{AI disclosure and nonpublication notice}

This manuscript was generated by an AI coding and mathematical reasoning system
from the repository materials cited below.  No human authorship, human review,
independent verification, peer review, acceptance, publication, or external
announcement is claimed.

\begin{thebibliography}{99}
\bibitem{EFGW}
C.~Elphick, M.~Farber, F.~Goldberg, and P.~Wocjan,
\emph{Conjectured bounds for the sum of squares of positive eigenvalues of a
graph}, Discrete Math. 339 (2016), 2215--2223.

\bibitem{AKMPZ}
S.~Akbari, H.~Kumar, B.~Mohar, S.~Pragada, and S.~Zhang,
\emph{Refinement of a conjecture on positive square energy of graphs},
arXiv:2506.07264, 2025.

\bibitem{LTZ}
Y.~Liu, Q.~Tang, and S.~Zhang,
\emph{The positive and negative square-energy conjecture},
arXiv:2607.18031, 2026.

\bibitem{ExpositionModules}
\emph{LaTeX exposition modules for the rank-six synthesis}, repository proof
note, 2026;
\path{positive-square-energy/hexacyclic-general/latex-exposition-modules.md}.

\bibitem{MultiblockAudit}
\emph{Hexacyclic multiblock theorem: exhaustive ledger and hostile audit},
repository proof note, 2026;
\path{positive-square-energy/hexacyclic-general/multiblock-items1-7-combined-theorem-audit.md}.

\bibitem{SingleBlockAudit}
\emph{Audit of the rank-six single-block implication}, repository proof note,
2026;
\path{positive-square-energy/hexacyclic-general/rank-six-single-block-implication-audit.md}.

\bibitem{MasterArchitecture}
\emph{Rank-six orders 2--10 master architecture}, promotion-gated repository
note, 2026;
\path{positive-square-energy/hexacyclic-general/rank-six-order2-10-master-architecture.md}.

\bibitem{OrderTenArchitecture}
\emph{Order-ten rank-six master architecture}, completion-gated repository
note, 2026;
\path{positive-square-energy/hexacyclic-general/order-ten-rank-six-master-architecture.md}.
\end{thebibliography}

\end{document}
